\section{\'Arbol de expansi\'on m\'inima (Kruskal)}
\whenuse{Conectar todos los v\'ertices con costo total m\'inimo. Kruskal es
c\'omodo con lista de aristas; Prim con adyacencias y grafos densos/impl\'icitos.}

Kruskal acepta pesos negativos. Si usa menos de $n-1$ aristas, el grafo es
disconexo y el resultado es un bosque de expansi\'on m\'inima.

\begin{lstlisting}[style=cpp]
struct MstEdge { int u, v; ll weight; };

pair<ll, vector<MstEdge>> kruskal(int n, vector<MstEdge> edges) {
    sort(edges.begin(), edges.end(), [](const MstEdge &a, const MstEdge &b) {
        return a.weight < b.weight;
    });
    DSU dsu(n);
    ll cost = 0;
    vector<MstEdge> chosen;
    for (const auto &edge : edges) if (dsu.unite(edge.u, edge.v)) {
        cost += edge.weight;
        chosen.push_back(edge);
    }
    return {cost, chosen};
}
\end{lstlisting}

\complexity{Kruskal $O(E\log E)$ tiempo y $O(V+E)$ memoria.}

Prim es conveniente cuando el grafo ya est\'a en listas de adyacencia.

\begin{lstlisting}[style=cpp]
optional<ll> prim(const vector<vector<pair<int,ll>>> &graph) {
    int n = graph.size();
    if (n == 0) return 0LL;
    vector<char> used(n, false);
    priority_queue<pair<ll,int>, vector<pair<ll,int>>,
                   greater<pair<ll,int>>> q;
    q.push({0, 0});
    ll cost = 0;
    int visited = 0;
    while (!q.empty()) {
        auto [weight, v] = q.top(); q.pop();
        if (used[v]) continue;
        used[v] = true; ++visited; cost += weight;
        for (auto [to, edge_weight] : graph[v])
            if (!used[to]) q.push({edge_weight, to});
    }
    if (visited != n) return nullopt;
    return cost;
}
\end{lstlisting}

\complexity{Prim $O(E\log V)$ tiempo y $O(V+E)$ memoria.}
